\documentclass[11pt]{article}
\usepackage[margin=1in]{geometry}
\usepackage[T1]{fontenc}
\usepackage{enumitem}
\setlength{\parindent}{0pt}
\setlength{\parskip}{6pt}

\begin{document}

\section*{2026-10-28\_read\_mmt-uncertainty}

\begin{tabular}{@{}p{0.14\textwidth}p{0.82\textwidth}@{}}
QUESTION & What does the uncertainty in the minimum mortality temperature do to the relative risks centred on it? \\[4pt]
READ & Lee, Kim, Hwang, Zanobetti, Schwartz and Chung 2017, ``Monte Carlo simulation-based estimation for the minimum mortality temperature in temperature-mortality association study'', main text and Additional file~1. \\[4pt]
FINDING & Centring on a fixed minimum barely moves the relative risk but breaks its interval. \\
\end{tabular}

\subsection*{Finding}

The minimum's uncertainty matters for the RR mainly through the width of its interval.
We went in thinking the tilt from centring on a noisy minimum was unmeasured and was the main harm.
Lee and colleagues measure it in simulation, and it is upward and small.
They planted four true curves and carried each simulated minimum into the RR.
In the U-shape the heat RR is high by 0.0005 on a true 1.079.
In the rotated S, where the true minimum sits just under the 99th percentile, it is high by 0.027 on a true 1.015 without a window, and by 0.003 with the 1st--99th window.
With the minimum fixed inside 1st--99th, the cold RR interval covers the truth 58.4\% of the time when the curve has a long flat floor.
The heat RR interval covers it 85.8\% of the time when the minimum sits near the ceiling.
Keeping only the simulated draws whose minimum falls inside 1st--99th brings heat coverage back to 98.1--100\%.
These are simulation results on one-dimensional moving-average curves, not on a 21-lag cross-basis.

\subsection*{Connections}

\begin{enumerate}[leftmargin=*]
\item The 1st--99th window fixes the point and costs the interval.
Tob\'ias 2017 introduced it to keep small-city minima off the tails.
England and Wales 2022 shows its ceiling binding, with heat RR at 1.000 in the top LSOA of every age group.
Lee 2017 prices it in simulation, where the same window drops coverage to 58.4\% on a floor and 85.8\% near the ceiling.
\item The window is chosen in two families and argued in neither.
The Lancet 2015, England and Wales 2022, Tob\'ias 2017 and Lee 2017 search 1st--99th.
Chen 2025 and \'Eric's Qu\'ebec script search 25th--90th.
\item Every paper that searches a fitted curve searches a shrunk curve, never the raw first-stage fit.
The Lancet 2015 and Lee 2017 search BLUPs, and England and Wales 2022 searches the downscaled LSOA curve.
\end{enumerate}

\subsection*{My work}

The minimum in \texttt{fit\_das} is found the way the field finds it, the lowest point inside 1st--99th, and nothing carries its uncertainty forward.
A CMA-level RR reported against that fixed minimum will have intervals that are too narrow wherever the curve has a flat floor or a minimum near the 99th.

\subsection*{Contradictions}

None.

\subsection*{Corrections}

Going in, the upward tilt was taken as the main harm and no RR in the set was recomputed with the minimum's uncertainty.
Lee 2017 recomputes the RR in every draw, shows the tilt is small, and shows the larger harm is interval coverage.

\subsection*{Questions}

\begin{enumerate}[leftmargin=*]
\item Do the 500 Monte Carlo draws behind England and Wales 2022's empirical intervals re-find the minimum in each draw, or hold it fixed?
Read the supplement's uncertainty annexe, or \texttt{02.fcastcomp.R} in Mistry's repository.
\item How far does coverage fall for a 21-lag cross-basis reduced to five coefficients, rather than a moving average?
Measure it in the outside simulation by adding the draw-and-refind step to \texttt{fit\_das} and scoring coverage against the planted p80 minimum.
\item Does Tobias et al.\ 2021 give an argument for 25th--90th that the 1st--99th papers lack?
\end{enumerate}

\subsection*{People}

Whanhee Lee and Yeonseung Chung designed the study and wrote it, per the contributions line.
Chung sits in mathematical sciences at KAIST.
Ho Kim, at Seoul National University's school of public health, appears on the Asian multi-city temperature papers the article cites.
Zanobetti and Schwartz at Harvard supplied the 135 US cities, the same data as the Lancet 2015 study.

\end{document}